% ============================================================================
% 第六节：镜像对称与递归
% ============================================================================

\section{镜像对称与递归结构}\label{sec:mirror}

镜像对称是计数几何最深刻的物理来源，
也是递归结构的核心应用场景。
本节阐述镜像对称如何通过微分方程与递归产生GW不变量。

\subsection{镜像对称的数学表述}

\begin{definition}[镜像对]\label{def:mirrorpair}
Calabi--Yau三维簇$X$与$X^\vee$构成\textbf{镜像对}，若：
\begin{enumerate}[label=(\arabic*)]
\item $X$与$X^\vee$的Hodge数满足$h^{p,q}(X) = h^{3-p,q}(X^\vee)$；
\item $X$的复结构形变对应$X^\vee$的Kähler形变（SYZ镜像对称）；
\item $X$的A-模型（GW理论）对应$X^\vee$的B-模型（周期积分）。
\end{enumerate}
\end{definition}

\subsection{Picard--Fuchs方程与$J$-函数}

\begin{definition}[Picard--Fuchs方程]\label{def:pf}
设$X^\vee \to \Delta$是镜像Calabi--Yau的复结构退化族，
$\Delta \subset \C$是圆盘。
周期积分$\Pi(t) = \int_{\gamma_t} \Omega_t$（$\Omega_t$是全纯3-形式）
满足\textbf{Picard--Fuchs方程}：
\begin{equation}\label{eq:pf}
\mathcal{L} \Pi(t) = 0, \quad
\mathcal{L} = \theta^N + \sum_{k=1}^N P_k(t) \theta^{N-k},
\end{equation}
其中$\theta = t \frac{d}{dt}$，$P_k$是多项式。
\end{definition}

\begin{definition}[$J$-函数]\label{def:jfunc}
亏格0的$J$-函数定义为
\begin{equation}\label{eq:jfunc}
J(q) = q \cdot \exp\left( \frac{F_0(q)}{H} \right),
\end{equation}
其中$F_0(q) = \sum_\beta \GWinv_{0,1,\beta} q^\beta$是亏格0势函数。
\end{definition}

\begin{theorem}[镜像定理，精确形式]\label{thm:mirror2}
$J$-函数等于Picard--Fuchs方程的特定解（$I$-函数）：
\[
J(q) = \frac{I(q)}{I_0(q)},
\]
其中$I_0$是$I$-函数的常数项。
\end{theorem}

\subsection{高亏格BCOV递归}

Bershadsky--Cecotti--Ooguri--Vafa（BCOV）\cite{BCOV1994}
提出了高亏格GW不变量的递归计算方案，
基于镜像Calabi--Yau上的holomorphic anomaly equation。

\begin{theorem}[BCOV holomorphic anomaly equation]\label{thm:bcov}
亏格$g$的GW势函数$F_g$满足
\begin{equation}\label{eq:bcov}
\bar\partial_i \partial_{\bar j} F_g
= \frac{1}{2} \ee^{2K} G_{i\bar j}^{a\bar b}
\left( \partial_a \partial_{\bar b} F_{g-1}
+ \sum_{g_1+g_2=g} \partial_a F_{g_1} \cdot \partial_{\bar b} F_{g_2} \right),
\end{equation}
其中$G_{i\bar j}^{a\bar b}$是Weil--Petersson度量的联络，
$K$是Kähler势。
\end{theorem}

BCOV方程\eqref{eq:bcov}与CEO递归\eqref{eq:toprec}具有相同的结构：
\textbf{亏格$g$的量由亏格$g-1$和更低亏格的组合生成}。
这一结构相似性不是巧合——
CEO递归可视为BCOV方程的普适化、几何无关的形式。

\subsection{局部镜像对称与精确解}

对局部Calabi--Yau（如$\Ocal(-1) \oplus \Ocal(-1) \to \CP^1$，即resolved conifold），
镜像对称可显式求解。

\begin{theorem}[Resolved conifold的全亏格GW不变量]\label{thm:conifold}
对resolved conifold $X = \Ocal(-1) \oplus \Ocal(-1) \to \CP^1$，
亏格$g$、度数$d$的GW不变量为
\begin{equation}\label{eq:conifold}
\GWinv_{g, d} = \frac{(-1)^{g-1} |B_{2g}|}{2g \cdot (2g-2)! \cdot d^{2g-1}},
\end{equation}
其中$B_{2g}$是Bernoulli数。
\end{theorem}

这一公式由Fab--Pandharipande通过拓扑顶点方法证明，
是高亏格GW理论的少数精确结果之一。
它也是MNOP对应下DT不变量计算的基准。

\subsection{Givental形式化与$R$-矩阵}

Givental将镜像定理重新表述为半单Frobenius流形上的$R$-矩阵作用，
这一形式化与CEO递归有深刻联系。

\begin{definition}[Givental $R$-矩阵]\label{def:rmatrix}
设$QH^*(X)$是半单量子上同调，
在基$\{\Delta_a\}$下$\Delta_a \star \Delta_b = \delta_{ab} \Delta_a$。
\textbf{$R$-矩阵}是形式级数
\begin{equation}\label{eq:rmatrix}
R(z) = \exp\left( \sum_{k=1}^\infty R_k z^{2k-1} \right),
\end{equation}
其中$R_k$是$H^*(X)$上的自同态。
\end{definition}

\begin{theorem}[Givental重建定理]\label{thm:givental}
亏格$g$、$n$-点的GW不变量由
\[
\GWinv_{g,n} = \frac{1}{n!} \int_{\Mbar_{g,n}}
\left( \sum_a \frac{T_a}{\sqrt{\Delta_a}} \right)^n
\cdot R(\psi) \cdot \text{topological terms}
\]
给出，其中$R(\psi)$是$R$-矩阵在$\psi$-类上的作用。
\end{theorem}

Givental的$R$-矩阵与CEO递归的核$K_\alpha$通过\textbf{Dubrovin--Givental对应}关联：
\textbf{$R$-矩阵的作用等价于在特定谱曲线上运行CEO递归}。
这一对应是拓扑递归统一性的最精确表述。
